% !Mode:: "TeX:UTF-8"
% 本部研究生与威海博士的开题报告正文，照研究生院《开题报告》模板“说明”里列的九项写。
% 括号里是模板的填写要求，写完删掉。第 9 项参考文献由模板排在最后，这里不写。
% \section、\subsection、\subsubsection 依次是节、条、款，编号与字号由模板给。

\section{课题来源及研究目的和意义}

\subsection{课题的来源}

本课题来源于□□□项目“□□□□□□”（项目编号：□□□□）。

\subsection{课题研究目的和意义}

（不少于 500 字）

发展国防工业、微电子工业等尖端技术需要精密和超精密的仪器设备，精密仪器设备要求高速、……

气体静压轴承由于具有运动精度高、摩擦损耗小、发热变形小、寿命长、无污染等特点，在航空航天工业、半导体工业、纺织工业和测量仪器中得到广泛应用。……

\section{国内外研究现状及分析}

\subsection{国外研究现状}

1828年，R.R.Willis\cite{zhao1998}发表了一篇关于小孔节流平板中压力分布的文章，这是有记载的研究气体润滑的最早文献。……

\subsection{国内研究现状}

（注意对所引用文献的准确标注）

国内对气体润滑的研究始于二十世纪六十年代\cite{lin1992,xin1999}，……

\subsection{国内外文献综述的简析}

（不少于 500 字；综合评述：国内外研究取得的成果，存在的不足或有待深入研究的问题）

……

\section{主要内容及方案}

（不少于 1500 字；撰写宜使用将来时态，切忌将论文目录直接作为主要内容，要突出本人研究内容）

本课题拟对局部多孔质气体静压止推轴承的静态特性和稳定性进行理论与实验研究，主要内容如下：

（1）建立多孔质石墨渗透率与分形维数之间关系的数学模型；

（2）……

图、表、公式照学位论文的写法，如\wref{tab:prop-1}。

\begin{table}[!ht]
  \centering
  \caption{1号试样渗透率测试数据}
  \label{tab:prop-1}
  \begin{booktabs}{colspec={cccc}, rows={m}}
    \toprule
    {供气压力\\ $P_{\mathrm{s}}$ (MPa)} & {流量测量\\ $M'$ (m$^3$/h)} & {压力差\\ $\Delta P$ (Pa)} & lg $M$ \\
    \midrule
    0.15 & 0.009 & 46 900 & -5.636 01 \\
    0.2 & 0.021 & 96 900 & -5.338 76 \\
    0.25 & 0.039 & 146 900 & -5.130 01 \\
    \bottomrule
  \end{booktabs}
\end{table}

\section{预期达到的目标}

……

\section{已完成的学术研究/实践工作}

（根据课题情况选择；详细撰写目前已进行的学术研究或者实践工作的内容和完成情况）

……

\section{进度安排}

（建议从进入课题时间开始）

……

\section{为完成课题所需的条件}

……

\section{研究过程中可能遇到的问题以及解决的措施}

……
